% Generated draft: requires mathematical and source review.
\documentclass[11pt]{article}
\usepackage[margin=25mm]{geometry}
\usepackage{fontspec}
\setmainfont{Latin Modern Roman}
\usepackage{amsmath,amssymb,mathtools,unicode-math}
\setmathfont{Latin Modern Math}
\usepackage{enumitem,needspace,xurl,hyperref}
\setlength{\parindent}{0pt}
\setlength{\parskip}{5pt}
\setlength{\emergencystretch}{3em}
\newcommand{\N}{\mathbb{N}}
\newcommand{\Z}{\mathbb{Z}}
\newcommand{\Q}{\mathbb{Q}}
\newcommand{\R}{\mathbb{R}}
\newcommand{\C}{\mathbb{C}}
\newcommand{\F}{\mathbb{F}}
\newcommand{\A}{\mathbb{A}}
\newcommand{\PP}{\mathbb{P}}
\newcommand{\E}{\mathbb{E}}
\newcommand{\eps}{\varepsilon}
\newcommand{\dd}{\,\mathrm{d}}
\begin{document}
\section*{UnsolvedMath ID: 30006531}
\textbf{Source problem code:} catalogue entry
\subsection{Definitions and mathematical statement}
An \emph{indexed forest} is a sequence \(F=(T_1,T_2,\ldots)\) of plane binary trees, with \(T_j\) a single leaf for all sufficiently large \(j\). Here the trees are interpreted as finite rooted trees in which every internal node has an ordered pair of children, designated left and right. Label all leaves consecutively by \(\Z_{>0}\), reading the leaves of each tree from left to right and reading the trees in sequence. Write \(I(F)\) for the set of internal nodes and \(\lambda(\ell)\) for the label of a leaf \(\ell\).

For a map \(\kappa:I(F)\to\Z_{>0}\), extend its values to leaves by setting \(\widetilde\kappa(\ell)=\lambda(\ell)\), and set \(\widetilde\kappa(v)=\kappa(v)\) for internal nodes. Call \(\kappa\) admissible if, for every internal node \(v\) with left and right children \(v_L,v_R\),
\[
\widetilde\kappa(v_L)\geq\kappa(v),\qquad
\widetilde\kappa(v_R)>\kappa(v).
\]
For commuting indeterminates \(x_1,x_2,\ldots\), define the forest polynomial
\[
P_F=\sum_{\kappa\text{ admissible}}\prod_{v\in I(F)}x_{\kappa(v)}
\in\Z[x_1,x_2,\ldots].
\]
The empty product is \(1\).

The explicit problem is to give a combinatorial or otherwise explicit description of the coefficients \(a^G_{i,F}\) in
\[
(x_1+\cdots+x_i)P_F=\sum_G a^G_{i,F}P_G,
\]
for every indexed forest \(F\) and every positive integer \(i\), with \(G\) ranging over indexed forests. The intended reading is a finite expansion as a polynomial identity. The supplied excerpt does not state the coefficient domain or the basis theorem that would guarantee existence and uniqueness of these coefficients; these matters require verification against the cited forest-polynomial source. The positive-integer range for \(i\) and the finite-tree convention are implicit interpretations of the supplied notation.

The broader request is to find forest counterparts of structural and positivity properties of Schubert polynomials. No particular list of such properties or precise positivity assertion is specified, so this broader request cannot be made into a single determinate mathematical statement from the supplied record alone.
\subsection{Short English statement}
Forest polynomials sum monomials over internal-node labelings constrained by the canonical leaf labels. Describe explicitly how multiplying any forest polynomial by \(x_1+\cdots+x_i\) expands in forest polynomials. The broader request for Schubert-property analogues remains unspecified, and the expansion conventions require source review.
\subsection{Sources}
\begin{itemize}
\item Local UnsolvedMath registry, \texttt{data/unsolvedmath/problems.json}, permanent ID \(30006531\), source problem code \texttt{catalogue entry}, title ``Forest Analogues of Schubert-Polynomial Properties.'' No individual problem proposer is identified in the supplied record.
\item Supplied primary provenance: \emph{Enumerative Combinatorics} (2026), Oberwolfach Reports, DOI \texttt{10.4171/owr/2026/2}, pp.~81--156. \url{https://doi.org/10.4171/owr/2026/2}. The supplied citation gives the report page range, without an exact problem page.
\item Supplied literature reference: Nadeau--Tewari, \emph{Advances in Mathematics}. \url{https://doi.org/10.1016/j.aim.2024.109834}. The supplied record does not give the article title or an exact theorem supporting the expansion conventions.
\item Supplied literature reference: Oberwolfach Report 2026/2. \url{https://doi.org/10.4171/owr/2026/2}.
\item Literature assessment dated exactly \texttt{2026-08-22}: supplied dated curation, not independently verified. It labels the problem open and states that neither the general forest-counterpart question nor the requested explicit coefficient description is settled; it also states that forest-polynomial theory supplies positive bases and expansions without providing the general Schubert-property counterpart requested here. The record explicitly requests independent verification of time-sensitive status and cited claims. Its \texttt{corrected\_verified} statement label and claim of purely typographical repairs are also supplied metadata, not independent mathematical verification.
\end{itemize}
\end{document}
